\documentclass[10pt,a4paper]{amsart}
\usepackage[margin=19mm]{geometry}
\usepackage{amsmath,amssymb,mathtools}
\usepackage[scr=rsfs,cal=euler]{mathalfa}
\usepackage{xfrac}
\usepackage[unicode,hidelinks,bookmarks=false]{hyperref}
\newcommand{\ie}{i.e.\ }
\newcommand{\cf}{cf.\ }
\newcommand{\resp}{respectively\ }
\newcommand{\Subsection}{\subsection}
\providecommand{\nobreakdash}{\nobreak\hbox{-}}
\setlength{\emergencystretch}{2em}
\multlinegap0pt
\usepackage{mathtools}
\usepackage{amssymb}
\usepackage{bbm}
\usepackage{xcolor}
\usepackage[all]{xy}
\usepackage{bm}
\usepackage[normalem]{ulem}
\usepackage[shortlabels]{enumitem}
\newtheorem{lemma}{Lemma}
\newcommand{\Gbf}{\mathbf{G}}
\newcommand{\Q}{\mathbb{Q}}
\newcommand{\Z}{\mathbb{Z}}
\title{Strengthenings of Mazur's Conjecture for Higher Heegner Points}
\author{Xiaoyu Zhang}
\date{Source: arXiv:2602.07210v1 (CC BY 4.0)}
\begin{document}
\maketitle

\noindent\emph{Source and licence.} Strengthenings of Mazur's Conjecture for Higher Heegner Points. Version arXiv:2602.07210v1 (CC BY 4.0), distributed by arXiv under the Creative Commons Attribution 4.0 International licence, http://creativecommons.org/licenses/by/4.0/, as recorded in the arXiv licence metadata for this exact version and shown on the abstract page https://arxiv.org/abs/2602.07210v1. Full licence notice: \texttt{licenses/arxiv-2602.07210-CC-BY-4.0.txt}.\par\smallskip

\noindent\textit{Editorial scope.} Counted unit: the article's lemma orbit (source0.tex lines 1906--1909). The article's complete original proof of it is delivered with it (source0.tex lines 1910--1974, one \texttt{proof} environment of 65 lines). No mathematics is written, repaired or re-derived: the statement and the proof are the article's own lines, in the article's own notation, and no retained line is substituted or re-flowed.  Retained uncounted as the source-local context the proof needs: lines 1816--1822.  Nothing was omitted from any retained span, so the package carries no omission record.  No retained line contains a citation, so the wrapper carries no bibliography and no bibliography entry.  No retained line contains a figure environment, an image-inclusion command or drawing code, so no image is delivered and the package declares no asset dependency.  Editorial formatting only, in addition: an AMS \texttt{amsart} preamble that restates the article's own declarations and macros used by the retained lines, namely the environments \texttt{lemma} on separate counters, together with the standard packages those lines need.  The retained environments print in this document as Lemma 1 in order of appearance, whereas the article prints the same items with the numbering of its own full document.  The complete native source of the article as distributed in the arXiv e-print for this exact version is bundled unchanged as \texttt{sources/194158/source0.tex}.
	\begin{equation}\label{G_0(Z_q)=R_q}
		\Gbf_0(\Z_q)
		\simeq
		G(\Z_q),
		\quad
		q\neq\ell.
	\end{equation}

		\begin{lemma}\label{Galois orbit contained in Hecke orbit}
			For any distinct $g,g'\in S_n$, the right cosets
			$gh[n]G(\Z_q)$ and $g'h[n]G(\Z_q)$ are distinct right cosets in the decomposition $G(\Z_q)a(n)G(\Z_q)$ (here we use the identifications (\ref{G_0(Z_q)=R_q})).
		\end{lemma}

		\begin{proof}
			We first show that the two cosets are distinct.
			If $q$ splits in $L$, write $g=\begin{pmatrix}
				a & 0 \\ 0 & 1
			\end{pmatrix}$ and $g'=\begin{pmatrix}
				a' & 0 \\ 0 & 1
			\end{pmatrix}$. Then we have
			\[
			gh[n]G(\Z_q)
			=
			\begin{pmatrix}
				a\varpi_q^e & a \\ 0 & 1
			\end{pmatrix}
			G(\Z_q)
			=
			\begin{pmatrix}
				n & a \\ 0 & 1
			\end{pmatrix}
			G(\Z_q),
			\quad
			g'h[n]G(\Z_q)
			=
			\begin{pmatrix}
				n & a' \\ 0 & 1
			\end{pmatrix}
			G(\Z_q).
			\]
			Since 
			\[
			\begin{pmatrix}
				n & a \\ 0 & 1
			\end{pmatrix}^{-1}
			\begin{pmatrix}
				n & a' \\ 0 & 1
			\end{pmatrix}
			=
			\begin{pmatrix}
				1 & (a'-a)/n \\ 0 & 1
			\end{pmatrix}\notin G(\Z_q),
			\]
			these right cosets are indeed distinct.
			
			
			
			
			If $q$ is inert in $L$, write
			$g^{-1}g'=\begin{pmatrix}
				a & b \\ c & d
			\end{pmatrix}$. By definition of $S_n$, one checks that $b/n\notin\Z_q$. Then we have
			\[
			(gh[n])^{-1}(g'h[n])
			=
			\begin{pmatrix}
				a & b/n \\ cn & d
			\end{pmatrix}\notin G(\Z_q).
			\]
			It follows again that these two right cosets are distinct.
			
			
			
			
			Next we show that both cosets lie in $G(\Z_q)a(n)G(\Z_q)$. By the Cartan decomposition for $G(\Q_q)$, we can characterize $G(\Z_q)a(n)G(\Z_q)$ as the subset of $G(\Q_q)$, consisting of matrices $A=\begin{pmatrix}
				a & b \\ c & d
			\end{pmatrix}$ with $a,b,c,d\in\Z_q$ such that $\det(A)\in n\Z_q^\times$ and $a,b,c,d$ are not simultaneous divisible by $q$. Using this criterion, it follows easily that $gh[n]G(\Z_q)$ and $g'h[n]G(\Z_q)$ both lie in $G(\Z_q)a(n)G(\Z_q)$.
		\end{proof}

\end{document}
